\section*{Chapitre \ref{ch:pl:simulation-production-parity} --- Parité simulation--production}

\begin{solution}{exo:pl:simulation-production-parity:1}
Ni 99.99 ni 0.01 n'ont de représentation binaire exacte, et leur différence est stockée sous la forme
\texttt{99.97999999999999}~; multipliée par dix mille, elle se situe juste
sous 999\,800, et \texttt{int()} tronque. Il faut conserver les prix sous forme de nombres entiers de la plus petite unité (999\,900 $-$ 100)~; si un nombre à virgule flottante doit être
converti, il faut arrondir à l'unité la plus proche, jamais tronquer.
\end{solution}

\begin{solution}{exo:pl:simulation-production-parity:2}
Les minuteries déclenchées en production sont des entrées consignées dans le journal, aux instants où elles se sont déclenchées. Les réarmer dans le rejeu
reviendrait à livrer chaque minuterie deux fois, ou à un autre instant si la stratégie calculait son délai autrement~; le rejeu doit reproduire ce que la
stratégie a reçu.
\end{solution}

\begin{solution}{exo:pl:simulation-production-parity:3}
En production, l'horloge de la machine est synchronisée sur celle de la plateforme, si bien que l'heure de la machine et celle de l'hôte concordent et que le défaut n'a aucun
effet visible là-bas~; il apparaît dès que les deux diffèrent --- dans toute simulation et tout rejeu.
\end{solution}

\begin{solution}{exo:pl:simulation-production-parity:4}
Chaque changement de cotation qui arrive avant les accusés de réception amène la stratégie à croire qu'elle n'a aucun ordre en cours et à en envoyer un autre, et chaque
nouvel ordre est lui-même un changement de cotation pour le tour suivant~; la tempête croît avec le taux de cotation pendant l'aller-retour de quarante microsecondes. La
passerelle d'ordres, sa limite d'exposition dans le pire cas l'arrête~: elle compte chaque ordre susceptible d'être encore exécuté et en a refusé 8\,406 parmi les envois.
\end{solution}

\begin{solution}{exo:pl:simulation-production-parity:5}
Lorsque la stratégie annulait un ordre dont l'exécution était survenue au même instant sans avoir encore été livrée, la passerelle a refusé l'annulation
localement, alors que l'adaptateur en processus l'envoyait~; la plateforme l'a rejetée comme trop tardive, et la stratégie a reçu un rejet que la production
n'a jamais envoyé, et le message supplémentaire de la plateforme a retardé de 500 nanosecondes le prochain ordre. La passerelle avait raison~: une annulation qui ne peut aboutir ne doit
pas atteindre la plateforme~; l'adaptateur en processus la refuse désormais aussi.
\end{solution}

\begin{solution}{exo:pl:simulation-production-parity:6}
Les deux comparaisons confrontent les sorties de la production à des deuxièmes exécutions différentes~: le simulateur négocie sa propre histoire après la divergence
(exécutions et cotations différentes), alors que le rejeu reçoit les entrées de la production et ne diffère que là où les décisions de la stratégie diffèrent.
Un défaut dont les effets se cumulent à travers le marché touche davantage de sorties dans la comparaison avec le simulateur.
\end{solution}

\begin{solution}{exo:pl:simulation-production-parity:7}
Non~: les premières différences dépendent de la session --- quand le premier prix exige une soustraction en virgule flottante qui tombe hors de la grille, quand une exécution et
un changement de cotation coïncident pour la première fois --- sauf pour l'horloge murale, dont la première différence est la minuterie à trente secondes, dans toute session où des cotations la précèdent, et
la sémantique des accusés de réception, qui diverge au premier changement de cotation suivant les premiers ordres.
\end{solution}

\begin{solution}{exo:pl:simulation-production-parity:8}
Le test compare le backtest à lui-même~: toute entrée que l'environnement de backtest fournit mal (une horloge, un type de prix, un ordre des événements, une
règle d'accusé de réception) est fournie de façon erronée dans les deux cas, et les ordres concordent. La parité doit comparer le backtest aux entrées de la production
(rejouer un journal de production) et un environnement de type production avec le simulateur sur la même session.
\end{solution}

\begin{solution}{pb:pl:simulation-production-parity:1}
\begin{enumerate}
\item L'horloge de l'hôte, les données de marché et la saisie d'ordres, par ses fonctions de rappel~; rien d'autre.
\item Afin que les entrées d'un même instant soient livrées dans l'ordre du \omterm{def:pl:backtest-engine-architecture:event}{modèle d'événements}, quel que soit l'ordre dans lequel l'environnement les a produites.
\item La passerelle du livre~13 (automate d'état, limiteur de débit, exposition), l'encodage et le découpage en trames sur un flux d'octets, et le journal des entrées.
\item Chaque entrée telle que livrée, avec son heure de réception, au format de journal des entrées du livre~13~: 8~octets pour l'heure et un enregistrement de 48~octets.
\item Les consigne sans les envoyer nulle part~: les sorties du rejeu sont comparées, et non exécutées.
\item Dans le simulateur, la minuterie ne se déclenche jamais dans l'heure et l'ouverture dure cinquante minutes~; dans le rejeu, la minuterie journalisée se déclenche, mais la
stratégie lit sa propre horloge de machine et recote largement, en se séparant à 30.1~secondes.
\item $99.99 - 0.01$ devient 999\,799, hors de la grille de cotation~; la plateforme le rejette.
\item Lorsqu'une exécution et un changement de cotation arrivent au même instant.
\item Si un ordre envoyé mais non acquitté compte comme ordre en cours~; la stratégie le renvoie à chaque changement de cotation pendant l'aller-retour.
\item Aucun des quatre ici~: chacun a aussi modifié ce que l'hôte du rejeu livrait ou lisait (son horloge, ses prix entiers, son ordre des événements, sa
sémantique)~; un défaut partagé par les environnements de backtest et de production passerait inaperçu.
\item 175 (horloge murale), 126 (prix à virgule flottante), 203 (ordre des rappels), 2 (sémantique des accusés de réception).
\item 539 sorties sur 553, 1\,686 sur 1\,690, 387 sur 463 et 8\,449 sur 8\,452 face au simulateur.
\item La sémantique des accusés de réception diverge à la deuxième entrée~; l'ordre des rappels diverge le plus tardivement, après 203.
\item Annulations d'ordres déjà terminés~: la passerelle les a refusées, l'adaptateur en processus les a envoyées.
\item 8\,887 sorties issues de 188\,079 entrées sur dix sessions, sans aucune différence dans l'une ou l'autre comparaison.
\item \textbf{Résultat nommé.} L'horloge murale diverge après 175 entrées, les prix à virgule flottante après 126, l'ordre des rappels après 203 et la sémantique
des accusés de réception diverge après 2, touchant 539, 1\,686, 387 et 8\,449 sorties dans l'heure~; après correction, les exécutions de type production, de simulateur et de rejeu concordent sur
chaque sortie des dix sessions.
\item Qu'elle ne lise que l'horloge et les entrées de l'hôte, et qu'elle passe le harnais face au journal de production et au simulateur sur la même
session.
\item Des prix entiers, une livraison selon le \omterm{def:pl:backtest-engine-architecture:event}{modèle d'événements}, et une seule sémantique de saisie d'ordres, testée par le harnais.
\item Deux environnements sur une même session détectent ce qu'un rejeu partage avec la production~; un rejeu du journal de production détecte ce que deux
simulateurs partagent.
\item Parce que l'hôte fixe l'heure, les prix, l'ordre et la sémantique pour chaque stratégie, et que le harnais les teste une fois pour toutes.
\end{enumerate}
\end{solution}

\begin{solution}{iq:pl:simulation-production-parity:1}
Au premier ordre qui diverge, et aux entrées qui le précèdent immédiatement, à partir d'un journal des entrées de la production rejoué dans le backtest~: ensuite à
ce que la stratégie lit sans que ce soit une entrée (horloges, aléa, ordre d'itération), et à la façon dont les environnements diffèrent par les types, l'ordre des événements
et la sémantique de la saisie d'ordres.
\iqlookfor{Journal, première différence, puis les coupables habituels.}
\end{solution}

\begin{solution}{iq:pl:simulation-production-parity:2}
Parce que la plupart des prix décimaux n'ont pas de représentation binaire exacte, de sorte que les calculs et les conversions produisent des valeurs hors de la grille de cotation, et les tests d'égalité
échouent~; les entiers exprimés dans la plus petite unité sont exacts, se comparent exactement et se hachent de façon identique dans tous les langages.
\iqlookfor{Erreur de représentation, troncature, ordres hors grille.}
\end{solution}

\begin{solution}{iq:pl:simulation-production-parity:3}
Lui donner une seule horloge injectée --- l'heure de l'entrée courante et les minuteries comme entrées --- et interdire toute autre source de temps~; journaliser ensuite les
déclenchements de minuteries en production et les rejouer, et tester avec un harnais dont l'horloge de la machine diffère de celle de l'hôte.
\iqlookfor{Horloge injectée, minuteries comme entrées.}
\end{solution}

\begin{solution}{iq:pl:simulation-production-parity:4}
Par une règle fixe que la plateforme applique partout --- celle du chapitre livre d'abord les données de marché, puis la saisie d'ordres, puis les minuteries --- car
sinon le même instant produit des décisions différentes selon les environnements~; la règle choisie compte moins que le fait qu'il n'y en ait qu'une.
\iqlookfor{Un ordre total, identique dans tous les environnements.}
\end{solution}

\begin{solution}{iq:pl:simulation-production-parity:5}
Un seul hôte et une seule interface de stratégie~; des adaptateurs pour le backtest, le simulateur et la production~; les entrées de la production journalisées~; un harnais qui
compare les sorties événement par événement au rejeu du journal et au simulateur sur la même session, signale la première
différence et ses entrées, et s'exécute en intégration continue avec des défauts introduits à dessein qu'il doit détecter.
\iqlookfor{Journal, deux comparaisons, première différence, défauts introduits à dessein comme tests du harnais.}
\end{solution}
